\section*{Hoofdstuk \ref{ch:g8:solids} --- Piramides en kegels}

\begin{solution}{exo:g8:solids:1}
Viervlak: $4$ \omterm{def:g5:solids:def}{zijvlakken}, $6$ \omterm{def:g5:solids:def}{ribben}, $4$ \omterm{def:g5:solids:def}{hoekpunten} ($4 + 4 = 6 + 2$).
Zeshoekig grondvlak: $7$ \omterm{def:g5:solids:def}{zijvlakken}, $12$ \omterm{def:g5:solids:def}{ribben}, $7$ \omterm{def:g5:solids:def}{hoekpunten}
($7 + 7 = 12 + 2$).
\end{solution}

\begin{solution}{exo:g8:solids:2}
$V = \frac13 \times 36 \times 10 = 120$ cm$^3$.

\omterm{def:g6:shapes:triangles}{Rechthoekig} grondvlak: $B = 24$ cm$^2$, dus
$V = \frac13 \times 24 \times 7.5 = 60$ cm$^3$.
\end{solution}

\begin{solution}{exo:g8:solids:3}
$V = \frac13 \pi \times 25 \times 9 = 75\pi \approx 236$ cm$^3$.
\end{solution}

\begin{solution}{exo:g8:solids:4}
\omterm{def:g7:areas:prism}{Cilinder}: $\pi \times 16 \times 12 = 192\pi \approx 603$ cm$^3$. \omterm{def:g8:solids:def}{Kegel}:
$\frac{192\pi}{3} = 64\pi \approx 201$ cm$^3$. Verhouding: de \omterm{def:g8:solids:def}{kegel} is een derde
van de \omterm{def:g7:areas:prism}{cilinder}.
\end{solution}

\begin{solution}{exo:g8:solids:5}
$56 = \frac13 \times B \times 8$, dus is $B = \frac{56 \times 3}{8} = 21$
cm$^2$.
\end{solution}

\begin{solution}{exo:g8:solids:6}
Het \omterm{def:g5:solids:def}{zijvlak} rolt uit tot een sector van een schijf waarvan de \omterm{def:g3:shapes:shapes}{radius} de
\emph{schuine hoogte} $s$ is (de afstand van de top tot de rand --- dat \omterm{def:g6:lines:objects}{lijnstuk}
ligt op het oppervlak), niet de hoogte $h$.
\end{solution}

\begin{solution}{exo:g8:solids:7}
$V = \frac13 \times 35^2 \times 22 = \frac13 \times 1225 \times 22
= \frac{26\,950}{3} \approx 8\,983$ m$^3$ --- ongeveer $9\,000$ kubieke meter.
\end{solution}

\begin{solution}{exo:g8:solids:8}
Hoogte: $h = \sqrt{13^2 - 5^2} = \sqrt{169 - 25} = \sqrt{144} = 12$ cm. \omterm{def:g2:measure:units}{Inhoud}:
$V = \frac13 \pi \times 25 \times 12 = 100\pi$ cm$^3$.
\end{solution}

\begin{solution}{exo:g8:solids:9}
Dezelfde \omterm{def:g3:shapes:shapes}{radius}, dus is de \omterm{def:g2:measure:units}{inhoud} van de \omterm{def:g8:solids:def}{kegel} een derde van die van de \omterm{def:g7:areas:prism}{cilinder}
\emph{bij dezelfde hoogte}: de vloeistof komt in de \omterm{def:g7:areas:prism}{cilinder} tot
$\frac{9}{3} = 3$ cm.
\end{solution}

\begin{solution}{exo:g8:solids:10}
Zand $=$ één \omterm{def:g8:solids:def}{kegel}: $V = \frac13 \pi \times 4 \times 4.5 = 6\pi \approx 18.8$
cm$^3$. De zandloper bevat twee zulke \omterm{def:g8:solids:def}{kegels}, dus vult het zand precies de \omterm{def:g3:division:half}{helft}
van de totale binneninhoud.
\end{solution}

\begin{solution}{exo:g8:solids:11}
\emph{1.} Ja: de formule $V = \frac13 Bh$ geldt voor elke stand van de top, zolang
$h$ de loodrechte afstand tot het vlak van het grondvlak is.
$V = \frac13 \times 36 \times 8 = 96$ cm$^3$.

\emph{2.} Diagonaal van het grondvlak: $\sqrt{6^2 + 6^2} = 6\sqrt2$ cm. De
langste zijribbe is de schuine zijde van een \omterm{def:g6:shapes:triangles}{rechthoekige driehoek} met
rechthoekszijden $6\sqrt2$ (in het vlak van het grondvlak) en $8$ (verticaal):
\[
\sqrt{(6\sqrt2)^2 + 8^2} = \sqrt{72 + 64} = \sqrt{136} \approx 11.7
\text{ cm}.
\]
\end{solution}

\begin{solution}{pb:g8:solids:1}
\textbf{1.} $G$ hoort bij het bovenvlak $EFGH$ en bij de twee \omterm{def:g5:solids:def}{zijvlakken} $BCGF$
en $DCGH$. De drie \omterm{def:g5:solids:def}{zijvlakken} die $G$ \emph{niet} bevatten, zijn het grondvlak
$ABCD$, het voorvlak $ABFE$ en het linkervlak $ADHE$ (ze hebben alle het
\omterm{def:g5:solids:def}{hoekpunt} $A$ gemeen, de hoek tegenover $G$). De drie \omterm{def:g8:solids:def}{piramides} zijn $ABCDG$,
$ABFEG$ en $ADHEG$: elk heeft een vierkant \omterm{def:g5:solids:def}{zijvlak} van de \omterm{def:g5:solids:def}{kubus} als grondvlak en
het verre \omterm{def:g5:solids:def}{hoekpunt} $G$ als top. Elk is een \omterm{def:g8:solids:def}{piramide} met vierkant grondvlak: $5$
\omterm{def:g5:solids:def}{zijvlakken}, $8$ \omterm{def:g5:solids:def}{ribben}, $5$ \omterm{def:g5:solids:def}{hoekpunten} (\cref{ex:g8:solids:count}).

\textbf{2.} De lange diagonaal $[AG]$ verbindt de enige twee \omterm{def:g5:solids:def}{hoekpunten} die bij
geen van de drie grondvlakken horen. Een draai over een derde slag om de \omterm{def:g6:lines:objects}{lijn}
$(AG)$ stuurt de \omterm{def:g5:solids:def}{kubus} naar zichzelf (de drie \omterm{def:g5:solids:def}{ribben} die uit $A$ vertrekken ---
naar $B$, $D$, $E$ --- worden cyclisch verwisseld, en zo ook de drie \omterm{def:g5:solids:def}{ribben} die in
$G$ aankomen). Ze brengt het grondvlak $ABCD$ naar $ADHE$, dan $ADHE$ naar
$ABFE$, en weer terug: de toppen blijven in $G$, de grondvlakken draaien rond.
Elke \omterm{def:g8:solids:def}{piramide} wordt zo op de volgende gebracht: de drie zijn identieke kopieën.

\textbf{3.} Drie identieke stukken die de \omterm{def:g5:solids:def}{kubus} vullen, verdelen zijn \omterm{def:g2:measure:units}{inhoud} in
gelijke delen:
\[
V_{\text{piramide}} = \frac{a^3}{3} .
\]

\textbf{4.} De \omterm{def:g8:solids:def}{piramide} $ABCDG$ heeft grondvlak $ABCD$, met \omterm{def:g6:measure:area}{oppervlakte} $a^2$, en
haar top $G$ ligt verticaal boven de hoek $C$, op hoogte $CG = a$ (een verticale
\omterm{def:g5:solids:def}{ribbe} van de \omterm{def:g5:solids:def}{kubus}). Zoals in \cref{exo:g8:solids:11} geldt de formule met de
loodrechte hoogte $h = a$:
\[
V = \frac13 \times a^2 \times a = \frac{a^3}{3} ,
\]
in volmaakte overeenstemming met vraag 3 --- de verdeling en de formule
bevestigen elkaar.

\textbf{5.} Voor $a = 6$: $V = \frac{6^3}{3} = \frac{216}{3} = 72$ cm$^3$ per
\omterm{def:g8:solids:def}{piramide}. De gietproef is deze verdeling in vloeibare vorm: het \omterm{def:g7:areas:prism}{prisma} op het
grondvlak $ABCD$ met hoogte $a$ is de \omterm{def:g5:solids:def}{kubus} zelf, en er gaan precies drie
piramidevullingen in --- hier zetten de drie vullingen zich zelfs meetkundig in
elkaar tot de \omterm{def:g5:solids:def}{kubus}.

\textbf{6.} Het middelpunt $O$ met de vier hoeken van elk \omterm{def:g5:solids:def}{zijvlak} verbinden geeft
zes \omterm{def:g8:solids:def}{piramides}, één per \omterm{def:g5:solids:def}{zijvlak}: vierkant grondvlak, top $O$. Het middelpunt ligt
ten opzichte van de zes \omterm{def:g5:solids:def}{zijvlakken} even goed geplaatst (elke draaiing of
spiegeling van de \omterm{def:g5:solids:def}{kubus} laat $O$ liggen en verwisselt de \omterm{def:g5:solids:def}{zijvlakken}), dus zijn de
zes \omterm{def:g8:solids:def}{piramides} identiek. De hoogte van elk is de afstand van $O$ tot een \omterm{def:g5:solids:def}{zijvlak}:
de \omterm{def:g3:division:half}{helft} van de zijde, $\frac{a}{2}$.

\textbf{7.} Zes identieke stukken verdelen de \omterm{def:g5:solids:def}{kubus}: elk $V = \frac{a^3}{6}$.

\textbf{8.} Met de formule: grondvlak van \omterm{def:g6:measure:area}{oppervlakte} $a^2$, hoogte $\frac a2$:
\[
V = \frac13 \times a^2 \times \frac{a}{2} = \frac{a^3}{6} .
\]
De twee antwoorden stemmen overeen.

\textbf{9.} Beide verdelingen gaan over \omterm{def:g8:solids:def}{piramides} met heel bijzondere
verhoudingen, uit een \omterm{def:g5:solids:def}{kubus} gesneden; een slanke \omterm{def:g8:solids:def}{piramide}, een scheve \omterm{def:g8:solids:def}{piramide} of
een \omterm{def:g8:solids:def}{kegel} kun je zo niet in elkaar zetten (geen drie \omterm{def:g8:solids:def}{kegels} vullen een \omterm{def:g7:areas:prism}{cilinder}).
Daarom wordt \cref{thm:g8:solids:volume} op dit niveau \emph{aangenomen}: de
verdelingen en de gietproef maken de $\frac13$ volkomen geloofwaardig, en het
eerlijke algemene bewijs --- met integralen --- staat in het bovenbouwvolume.

\textbf{10.} $V = \frac{6^3}{6} = 36$ cm$^3$; en met de formule:
$V = \frac13 \times 36 \times 3 = 36$ cm$^3$. Ze stemmen overeen.

\textbf{11.} In de driehoek $SAB$ worden de middens van de zijden $[SA]$ en
$[SB]$ verbonden door een \omterm{def:g6:lines:objects}{lijnstuk} dat \omterm{def:g6:lines:perp}{parallel} is met $(AB)$ en lengte
$\frac{AB}{2} = \frac{c}{2}$ heeft (\cref{thm:g8:midpoints:direct}). Hetzelfde
geldt op elk \omterm{def:g5:solids:def}{zijvlak}, dus vormen de vier middens een vierkant met zijde
$\frac c2$, met zijden \omterm{def:g6:lines:perp}{parallel} met het grondvlak. Op de verticale \omterm{def:g6:lines:objects}{lijn} door de
top gaat de doorsnede door het midden van de hoogte (weer de
middenparallelstelling, in een driehoek door $S$, het middelpunt van het
grondvlak en een \omterm{def:g5:solids:def}{hoekpunt} van het grondvlak): het snijvlak zit op hoogte
$\frac h2$.

\textbf{12.} Het bovenste stuk is een regelmatige \omterm{def:g8:solids:def}{piramide} met vierkant grondvlak
met zijde $\frac c2$ en hoogte $\frac h2$ (haar grondvlak is de doorsnede, haar
top is $S$). Haar \omterm{def:g2:measure:units}{inhoud} is
\[
\frac13 \times \left(\frac c2\right)^2 \times \frac h2
= \frac13 \times \frac{c^2}{4} \times \frac h2
= \frac{1}{8} \times \frac{c^2 h}{3}
= \frac{V}{8} :
\]
een achtste van het geheel --- niet de \omterm{def:g3:division:half}{helft}.

\textbf{13.} De afgeknotte \omterm{def:g8:solids:def}{piramide} houdt de \omterm{def:g3:division:remainder}{rest}: $V - \frac V8 = \frac{7V}{8}$.
De doorsnede op halve hoogte verdeelt de \omterm{def:g2:measure:units}{inhoud} in de verhouding $1 : 7$ --- het
bescheiden uitziende onderstuk bevat zeven keer de \omterm{def:g2:measure:units}{inhoud} van de spitse top.

\textbf{14.} Onder halve hoogte zit $\frac78$ van de \omterm{def:g2:measure:units}{inhoud}. Met
$V = \frac13 \times 35^2 \times 22 = \frac{26\,950}{3} \approx 8\,983$ m$^3$
(\cref{exo:g8:solids:7}):
\[
\frac78 \times \frac{26\,950}{3} \approx 7\,860 \approx
7\,900 \text{ m}^3
\]
op honderd nauwkeurig --- de campagne van de ``onderste \omterm{def:g3:division:half}{helft}'' poetst bijna acht
keer zo veel glasinhoud als die van de bovenste.

\textbf{15.} Elke afmeting verdubbelen vermenigvuldigt de \omterm{def:g6:measure:area}{oppervlakte} van het
grondvlak met $2^2 = 4$ en de hoogte met $2$, en dus de \omterm{def:g2:measure:units}{inhoud} met
$4 \times 2 = 8 = 2^3$. Vergroten met een factor $k$ vermenigvuldigt
\omterm{def:g6:measure:area}{oppervlakten} met $k^2$ en nog één lengte met $k$: \omterm{def:g2:measure:units}{inhouden} groeien met de factor
$k^3$. Het bovenste stuk van vraag 12 is een verkleinde kopie van de hele
\omterm{def:g8:solids:def}{piramide} met $k = \frac12$, dus is zijn \omterm{def:g2:measure:units}{inhoud}
$\left(\frac12\right)^3 = \frac18$ van het geheel --- de verklaring in één
regel.
\end{solution}
